% CP-VI-0040
% target_chapter: VI/23 — Fiber Products
% source_unit_id: IMP-DL-D27460A55B-P009
% source: Downloads/theory-of-algebraic-geometry-5.tex L6151-L6378
% solution_provenance: SOURCE_IMPLICIT_SOLUTION

\subsection*{Problem 10. Cartesian Diagrams}

\textbf{Problem.}
A commutative diagram
\[
\begin{tikzcd}
	A \arrow[r] \arrow[d] & B \arrow[d] \\
	C \arrow[r] & D
\end{tikzcd}
\]
is Cartesian if the induced morphism
\[
A\longrightarrow B\times_D C
\]
is an isomorphism.

Show that the following diagrams are Cartesian in any category.

\subsection{First diagram}

We are given morphisms
\[
X\longrightarrow S\longrightarrow T,
\qquad
Y\longrightarrow S\longrightarrow T.
\]
Then we have fibre products
\[
X\times_S Y
\qquad\text{and}\qquad
X\times_T Y.
\]
Also, the diagonal morphism
\[
\Delta:S\longrightarrow S\times_T S
\]
is induced by the two identity maps \(S\to S\).

The diagram is
\[
\begin{tikzcd}
	X\times_S Y \arrow[r] \arrow[d]
	&
	S \arrow[d,"\Delta"]
	\\
	X\times_T Y \arrow[r]
	&
	S\times_T S.
\end{tikzcd}
\]

We prove that it is Cartesian.

Let \(Q\) be any object. A morphism
\[
Q\longrightarrow
(X\times_T Y)\times_{S\times_T S}S
\]
is the same as the following data:

\[
x:Q\to X,
\qquad
y:Q\to Y,
\qquad
s:Q\to S,
\]
such that:

\[
x \text{ and } y \text{ agree after mapping to } T,
\]
and
\[
\Delta\circ s
=
(x_S,y_S),
\]
where
\[
x_S:Q\to X\to S,
\qquad
y_S:Q\to Y\to S.
\]

But the condition
\[
\Delta\circ s=(x_S,y_S)
\]
means precisely that
\[
s=x_S=y_S.
\]
Therefore the data above is exactly the same as a pair of morphisms
\[
x:Q\to X,
\qquad
y:Q\to Y
\]
such that
\[
x_S=y_S.
\]
By the universal property of the fibre product, this is exactly the same
as a morphism
\[
Q\longrightarrow X\times_S Y.
\]

Thus for every object \(Q\), there is a natural bijection
\[
\operatorname{Hom}
\left(Q,X\times_S Y\right)
\cong
\operatorname{Hom}
\left(
Q,(X\times_T Y)\times_{S\times_T S}S
\right).
\]
By the Yoneda principle, this gives an isomorphism
\[
X\times_S Y
\cong
(X\times_T Y)\times_{S\times_T S}S.
\]
Therefore the diagram is Cartesian.

\subsection{Second diagram}

Now assume \(X\) and \(Y\) are objects over \(S\), and let
\[
f:X\longrightarrow Y
\]
be a morphism over \(S\).

The diagonal morphism is
\[
\Delta:Y\longrightarrow Y\times_S Y.
\]

The graph morphism of \(f\) is
\[
\Gamma_f:X\longrightarrow X\times_S Y,
\]
defined by
\[
\Gamma_f=(\operatorname{id}_X,f).
\]

The diagram is
\[
\begin{tikzcd}
	X \arrow[r,"f"] \arrow[d,"\Gamma_f"']
	&
	Y \arrow[d,"\Delta"]
	\\
	X\times_S Y \arrow[r,"f\times \operatorname{id}_Y"']
	&
	Y\times_S Y.
\end{tikzcd}
\]

We prove that it is Cartesian.

Let \(Q\) be any object. A morphism
\[
Q\longrightarrow
(X\times_S Y)\times_{Y\times_S Y}Y
\]
is the same as the following data:

\[
x:Q\to X,
\qquad
y:Q\to Y,
\qquad
z:Q\to Y,
\]
such that \((x,y)\) defines a morphism
\[
Q\to X\times_S Y,
\]
and
\[
(f\circ x,\ y)=\Delta\circ z.
\]
But
\[
\Delta\circ z=(z,z).
\]
Hence the condition becomes
\[
(f\circ x,\ y)=(z,z).
\]
Therefore
\[
z=f\circ x
\qquad\text{and}\qquad
y=f\circ x.
\]
Thus the whole data is uniquely determined by the single morphism
\[
x:Q\to X.
\]

Conversely, any morphism \(x:Q\to X\) gives such a triple by setting
\[
y=f\circ x,
\qquad
z=f\circ x.
\]
Therefore there is a natural bijection
\[
\operatorname{Hom}(Q,X)
\cong
\operatorname{Hom}
\left(
Q,(X\times_S Y)\times_{Y\times_S Y}Y
\right).
\]
By the Yoneda principle,
\[
X
\cong
(X\times_S Y)\times_{Y\times_S Y}Y.
\]
Therefore the second diagram is Cartesian.
